\documentclass{article}


% 中文支持
\usepackage[UTF8,heading = true]{ctex}


% 页面布局
\usepackage[a4paper]{geometry}
\usepackage{fancyhdr}
\pagestyle{fancy}
\newcommand{\Lhead}[1]{\fancyhead[L]{#1}}
\newcommand{\Chead}[1]{\fancyhead[C]{#1}}
\newcommand{\Rhead}[1]{\fancyhead[R]{#1}}
\newcommand{\Lfoot}[1]{\fancyfoot[L]{#1}}
\newcommand{\Cfoot}[1]{\fancyfoot[C]{#1}}
\newcommand{\Rfoot}[1]{\fancyfoot[R]{#1}}
\renewcommand{\headrulewidth}{0.4pt}
\renewcommand{\footrulewidth}{0pt}
\usepackage{lastpage}
%待填信息
% -----------------------------------------------------
\newcommand{\Course}{深度学习平台与应用}             		%      |
\newcommand{\Assignment}{深度学习作业 4}		  %      | 
\newcommand{\Name}{田永铭}				    	 %      |
\newcommand{\ID}{221900180}				  	   %       |
\newcommand{\Email}{tianym2022@smail.nju.edu.cn}  %    |
% ------------------------------------------------------
\fancyhead[L]{\Course}
\fancyhead[C]{\Assignment}
\fancyhead[R]{\ID \Name}
\fancyfoot[C]{\thepage\ / \pageref*{LastPage}}	
\setlength\headheight{36pt}
% 自定义页面样式第一页只有页脚没有页眉
\makeatletter
\def\ps@onlyfoot{
	\let\@oddhead\@empty
	\let\@evenhead\@empty
	\def\@oddfoot{\hfil\thepage\ / \pageref*{LastPage}\hfil}
	\def\@evenfoot{\hfil\thepage\ / \pageref*{LastPage}\hfil}
}
\makeatother
% maketitle
\title{\Huge\Assignment}
\author{\ID \Name \quad\href{mailto:tianym2022@smail.nju.edu.cn}{\color{magenta}\Email}}
\date{\today}


% 字体和颜色
\usepackage{fontspec} 
\fontsize{12pt}{\baselineskip}\selectfont
\usepackage{xcolor}


% 数学支持
\usepackage{amssymb}
\usepackage{amsfonts}
\usepackage{amsmath}
\usepackage{esint}
\usepackage{gensymb}
\usepackage{mathtools}
\usepackage{amsthm}

% 左边大括号
%\[
%\left\{
%\begin{aligned}
%	& \text{approximate  covariance matrix of GP} \\
%	& \text{only use subset}\\
%	& \text{conditional independence assumption}
%\end{aligned}
%\right.
%\]

% 矩阵
%\[
%\begin{matrix}
%	a & b \\
%	c & d
%\end{matrix}
%\]

% 行列式
%\[
%\begin{vmatrix}
%	a & b \\
%	c & d
%\end{vmatrix}
%\]


% 标题格式
\usepackage{titlesec}
% 自定义 section 标题格式
\renewcommand{\thesection}{\chinese{section}、}
\titleformat{\section}
{\normalfont\Large\bfseries\centering\color[rgb]{0,0,0.8}} % 设定字体和大小
{\thesection}{0em} % 设定编号和标题之间的距离
{}
% 自定义 subsection 标题格式
\renewcommand{\thesubsection}{\arabic{section}.\arabic{subsection} }
\titleformat{\subsection}
{\normalfont\large\bfseries\color[rgb]{0.23,0.65,0.99}}
{\thesubsection}{0em}
{}
% 自定义 subsubsection 标题格式
\renewcommand{\thesubsubsection}{\arabic{section}.\arabic{subsection}.\arabic{subsubsection} }
\titleformat{\subsubsection}
{\normalfont\normalsize\bfseries}
{\thesubsubsection}{0em}
{}
% 设置 section 标题的间距
\titlespacing*{\section}
{0pt}{3.5ex plus 1ex minus .2ex}{2.3ex plus .2ex}
% 设置 subsection 标题的间距
\titlespacing*{\subsection}
{0pt}{3.25ex plus 1ex minus .2ex}{1.5ex plus .2ex}
% 设置 subsubsection 标题的间距
\titlespacing*{\subsubsection}
{0pt}{3.25ex plus 1ex minus .2ex}{1.5ex plus .2ex}


% 代码 
% 如非默认，需使用[language=xxx,style=mystyle]
\usepackage{listings}
% language:  {[x86masm]Assembler}(汇编); python; C; C++; Java; SQL; Matlab; R; bash; make(makefile)
\lstdefinestyle{mystyle}{
	basicstyle=\ttfamily\small,
	numbers=left,
	numberstyle=\tiny\color{magenta},
	stepnumber=1,
	numbersep=10pt,
	backgroundcolor=\color[rgb]{0.95,0.95,0.95},
	showspaces=false,
	showstringspaces=false,
	showtabs=false,
	frame=single,
	rulecolor=\color{black},
	tabsize=2,
	captionpos=b,
	breaklines=true,
	breakatwhitespace=false,
	title=\lstname,
	keywordstyle=\color{blue},
	commentstyle=\color[rgb]{0.0,0.5,0.0},
	stringstyle=\color[rgb]{0.23,0.65,0.99}\itshape,
	escapeinside={\%*}{*)},
	morekeywords={}, % 添加更多关键字
	literate=
	{[}{{{\color{red}[}}}1
	{]}{{{\color{red}]}}}1
	{\{}{{{\color[rgb]{0.9,0,0.9}\{}}}1
	{\}}{{{\color[rgb]{0.9,0,0.9}\}}}}1
%	{0}{{{\color[rgb]{0.2,0.9,0.4}0}}}1
%	{1}{{{\color[rgb]{0.2,0.9,0.4}1}}}1
%	{2}{{{\color[rgb]{0.2,0.9,0.4}2}}}1
%	{3}{{{\color[rgb]{0.2,0.9,0.4}3}}}1
%	{4}{{{\color[rgb]{0.2,0.9,0.4}4}}}1
%	{5}{{{\color[rgb]{0.2,0.9,0.4}5}}}1
%	{6}{{{\color[rgb]{0.2,0.9,0.4}6}}}1
%	{7}{{{\color[rgb]{0.2,0.9,0.4}7}}}1
%	{8}{{{\color[rgb]{0.2,0.9,0.4}8}}}1
%	{9}{{{\color[rgb]{0.2,0.9,0.4}9}}}1
}
\lstset{style=mystyle}


% 图表
\usepackage{graphicx}
\usepackage{booktabs}
\usepackage{caption}
\captionsetup{font={small}}
\usepackage{diagbox}
\usepackage{array}
\usepackage{pgfplots}
\pgfplotsset{compat=1.17}

% 表1
%\begin{table}[h!]
%	\centering
%	\caption{Performance comparison on mini-ImageNet and tiered-ImageNet datasets.}
%	\begin{tabular}{@{}lcccccccc@{}}
%		\toprule
%		Dataset & Method & Backbone & 5-way 1-shot \\
%		\midrule
%		\multirow{9}{*}{\rotatebox[origin=c]{90}{mini-ImageNet}} 
%		& ProtoNet & Conv-4 & 46.291 \\
%		& MAML & Conv-4 & 49.704 \\
%		& BOIL & Conv-4 & 48.456 \\
%		& CAN & ResNet-12 & 59.493\\
%		& LEO & ResNet-12 & 55.056 \\
%		& SKD & ResNet-12 & 66.700  \\
%		& R2D2 & ResNet-12 & 59.802 \\
%		& DiffKendall (Ours) & ResNet-12 & \textbf{68.727} \\
%		\midrule
%		\multirow{6}{*}{\rotatebox[origin=c]{90}{tiered-ImageNet}} 
%		& ProtoNet & Conv-4 & 46.653 \\
%		& Relation Networks & Conv-4 & 55.107  \\
%		& MAML & Conv-4 & 51.038 \\
%		& R2D2 & ResNet-12 & 64.916  \\
%		& DiffKendall (Ours) & ResNet-12 & \textbf{66.598} \\
%		\bottomrule
%	\end{tabular}
%\end{table}

% 表2
%\begin{table}[t]
%	\caption{Classification accuracies for naive Bayes and flexible
%		Bayes on various data sets.}
%	\label{sample-table}
%	\vskip 0.15in
%	\begin{center}
%		\begin{small}
%			\begin{sc}
%				\begin{tabular}{lcccr}
%					\toprule
%					Data set & Naive & Flexible & Better? \\
%					\midrule
%					Breast    & 95.9$\pm$ 0.2& 96.7$\pm$ 0.2& $\surd$ \\
%					Cleveland & 83.3$\pm$ 0.6& 80.0$\pm$ 0.6& $\times$\\
%					Glass2    & 61.9$\pm$ 1.4& 83.8$\pm$ 0.7& $\surd$ \\
%					Credit    & 74.8$\pm$ 0.5& 78.3$\pm$ 0.6&         \\
%					Horse     & 73.3$\pm$ 0.9& 69.7$\pm$ 1.0& $\times$\\
%					Meta      & 67.1$\pm$ 0.6& 76.5$\pm$ 0.5& $\surd$ \\
%					Pima      & 75.1$\pm$ 0.6& 73.9$\pm$ 0.5&         \\
%					Vehicle   & 44.9$\pm$ 0.6& 61.5$\pm$ 0.4& $\surd$ \\
%					\bottomrule
%				\end{tabular}
%			\end{sc}
%		\end{small}
%	\end{center}
%	\vskip -0.1in
%\end{table}


% 算法
%\usepackage{algorithm}
%\usepackage{algorithmic}
%\begin{algorithm}[tb]
%	\caption{Bubble Sort}
%	\label{alg:example}
%	\begin{algorithmic}
	%		\STATE {\bfseries Input:} data $x_i$, size $m$
	%		\REPEAT
	%		\STATE Initialize $noChange = true$.
	%		\FOR{$i=1$ {\bfseries to} $m-1$}
	%		\IF{$x_i > x_{i+1}$}
	%		\STATE Swap $x_i$ and $x_{i+1}$
	%		\STATE $noChange = false$
	%		\ENDIF
	%		\ENDFOR
	%		\UNTIL{$noChange$ is $true$}
	%	\end{algorithmic}
%\end{algorithm}

% 超链接
\usepackage{hyperref}


% 论文参考文献
% 需在正文结尾加入 \bibliography{references}，在同目录下创建references.bib
\bibliographystyle{plain}


% 表情
% https://ctan.math.washington.edu/tex-archive/graphics/pgf/contrib/tikzsymbols/tikzsymbols.pdf
\usepackage{tikzsymbols}
\usepackage{tikz}
% 背景颜色：\colorbox{yellow}{...}
% 3D版本： 在前面加上d,比如：dSmiley

% 常用
% 笑脸：\Smiley[⟨scale⟩][⟨color⟩]
% 哭脸：\Sadey[⟨scale⟩][⟨color⟩]
% 生气脸： \Annoey[⟨scale⟩][⟨color⟩]
% 大笑： \Laughey[⟨scale⟩][⟨color⟩][⟨mouth color⟩]
% 挤眉弄眼： \Winkey[⟨scale⟩][⟨color⟩]
% 去世： \Xey[⟨scale⟩][⟨color⟩]
% 升天： \wInnocey[⟨scale⟩]
% 口吐芬芳： \Vomey[⟨scale⟩][⟨color⟩][⟨vomit color⟩]
% 猫： \Cat[⟨scale⟩]
% 炸弹： \Ninja[⟨scale⟩][⟨color⟩][⟨headband color⟩][⟨eye color⟩]
% 困： \Sleepey[⟨scale⟩][⟨color⟩][⟨cap color⟩][⟨star color⟩]
% 口罩：\Maskey[⟨scale⟩][⟨color⟩][⟨mask color⟩]
% 渐变弧度脸(-2到2)： \Changey[⟨scale⟩][⟨color⟩]{⟨mood⟩}
% 树（春夏秋冬）： \Springtree[⟨scale⟩]
% 咖啡： \Coffeecup[⟨scale⟩]



% 简化命令
\newcommand{\hs}[1]{\hspace{{#1}em}} %空格
\newcommand{\qd}{\quad} %空格
\newcommand{\qqd}{\qquad} %空格
% 链接
\newcommand{\myhref}[2]{\href{#1}{{\color[rgb]{0.9,0,0.9}{#2}}}}
\newcommand{\myred}[1]{{\color{red}{#1}}} %红色
% 图 args：figures下图片路径，大小，图片显示名
\newcommand{\myfig}[3]{
	\begin{figure}[h!]
		\centering
		\includegraphics[width=#2\linewidth]{figures/#1}
		\caption{#3}
		\label{fig:#1}
	\end{figure}
}
% 左右分栏插入图片
% args: 左图文件名，左图宽度，左图标题，右图文件名，右图宽度，右图标题
\newcommand{\mytwocolumnfig}[6]{
	\begin{figure}[h!]
		\centering
		\begin{minipage}[h!]{0.45\linewidth}
			\centering
			\includegraphics[width=#2\linewidth]{figures/#1}
			\caption{#3}
			\label{fig:#1}
		\end{minipage}
		\hfill
		\begin{minipage}[h!]{0.45\linewidth}
			\centering
			\includegraphics[width=#5\linewidth]{figures/#4}
			\caption{#6}
			\label{fig:#4}
		\end{minipage}
	\end{figure}
}


\newtheoremstyle{Problem}
{}
{}
{}
{}
{\bfseries}
{}
{ }
{Problem\thmnumber{ #2}\thmname{ #1}\thmnote{ (#3)} }
\theoremstyle{Problem}

\newtheorem{prob}{}
\newtheoremstyle{Solution}
{}
{}
{}
{}
{\bfseries}
{:}
{ }
{\thmname{#1}}
\makeatletter
\makeatother
\theoremstyle{Solution}
\newtheorem*{sol}{Solution}


\begin{document}
	\maketitle
	\begin{prob}
		\textbf{损失函数计算}\\
		已知一个线性分类器的预测得分为以下矩阵：
		$
		S = \begin{bmatrix}
			4.8&2.5&2.0\\
			-3.1&0&12.9\\
		\end{bmatrix}
		$
		
		(a) 其中每一行表示一个样本的预测得分，每一列表示一个类别。假设正确类别的索引为[0,2],请计算SVM的损失函数值，公式如下：
		$
		L = \frac{1}{N} \sum_{i=1}^{N} \sum_{j \neq y_i} \max(0, s_j - s_{y_i} + 1)
		$
		
		
		(b) 对于这个线性分类器的预测，计算出softmax概率以及交叉熵损失。
	\end{prob}
	
	\begin{sol}
		解：\\
	(a)  \begin{align*}
		L &= \frac{1}{2}\sum_{i=1}^{2}\sum_{j\neq y_i}max(0,s_j-s_{y_i}+1)\\
		&= \frac{max(0,2.5-4.8+1)+max(0,2.0-4.8+1)}{2} + \frac{max(0,-3.1-12.9+1)+max(0,0-12.9+1)}{2}\\
		&= 0\ .
	\end{align*}
	所以SVM损失函数值为0.\\
	(b) \begin{itemize}
		\item \textbf{计算softmax概率：}\\
		根据公式:
		\[P(y=k|x_i) = \frac{e^{s_k}}{\sum_{j=0}^2 e^{s_j}},\]得:(\myred{注意我这里都保留了两位小数})
		
		S\_{softmax} = $\begin{bmatrix}
		0.86	& 0.09&0.05 \\
		0.00	&0.00 &1.00
		\end{bmatrix}$
		\item \textbf{计算交叉熵损失：}\\
		根据单样本的交叉熵损失公式$L_i = -log\frac{e^{y_i}}{\sum_j e^{s_j}}$，得：(\myred{注意我这里都保留了两位小数})
		$L_{\text{entropy}} = -\frac{\log\frac{0.86}{0.86+0.09+0.05}+\log\frac{1.00}{1.00+0.00+0.00}}{2} \approx 0.07\ .$
	\end{itemize}
	\end{sol}
	
	
	\begin{prob}
		\textbf{numpy 实现带L2正则化的逻辑回归}\\
		\begin{lstlisting}[language=python,style=mystyle]
import numpy as np
import matplotlib.pyplot as plt


# 生成高斯分布的二分类数据
	def generate_data(n_samples=100,  mean1=(2, 2), mean2=(-2, -2), cov=((1, 0), (0, 1))):
	np.random.seed(42) 
	X1 = np.random.multivariate_normal(mean1, cov, n_samples)  # 正类
	X2 = np.random.multivariate_normal(mean2, cov, n_samples)  # 负类
	Y1 = np.ones((n_samples, 1))  
	Y2 = np.zeros((n_samples, 1))  
	X = np.vstack((X1, X2))  
	Y = np.vstack((Y1, Y2))  
	indices = np.arange(X.shape[0])
	np.random.shuffle(indices) 
	return X[indices], Y[indices]


# Sigmoid 函数
def sigmoid(z):
	return 1 / (1 + np.exp(-z))


# 逻辑回归（带 L2 正则化）
def logistic_regression(X, Y, lr=0.01, epochs=1000, lambda_=0.1):
	m, n = X.shape  # 样本数量和特征数量
	W = np.zeros((n, 1))  # 初始化权重
	b = 0  # 初始偏置
	losses = []

	for epoch in range(epochs):
		# 前向传播
		Z = np.dot(X, W) + b  # 线性部分
		A = sigmoid(Z)  # Sigmoid 激活
		#todo 损失函数（交叉熵 + L2 正则化项）
		#loss = ?
		losses.append(loss)
		
		#todo 反向传播
		#dW = ?  交叉熵 + L2 正则化 梯度
		#db = ? 偏置的梯度
		
		# 参数更新
		W -= lr * dW
		b -= lr * db
	
	return W, b, losses


# 预测函数
def predict(X, W, b):
	Z = np.dot(X, W) + b
	A = sigmoid(Z)  # 激活函数
	return (A >= 0.5).astype(int)


# 绘制数据分布和决策边界
def plot_decision_boundary(X, Y, W, b):
	plt.figure(figsize=(8, 6))
	plt.scatter(X[Y.flatten() == 0][:, 0], X[Y.flatten() == 0][:, 1], color='red', label='Class 0')
	plt.scatter(X[Y.flatten() == 1][:, 0], X[Y.flatten() == 1][:, 1], color='blue', label='Class 1')
	x_min, x_max = X[:, 0].min() - 1, X[:, 0].max() + 1
	y_min, y_max = X[:, 1].min() - 1, X[:, 1].max() + 1
	xx, yy = np.meshgrid(np.arange(x_min, x_max, 0.01), np.arange(y_min, y_max, 0.01))
	grid = np.c_[xx.ravel(), yy.ravel()]
	Z = predict(grid, W, b)
	Z = Z.reshape(xx.shape)
	plt.contourf(xx, yy, Z, alpha=0.3, cmap="coolwarm")
	plt.legend()
	plt.title("Decision Boundary")
	plt.show()


# 主程序
if __name__ == "__main__":
	X, Y = generate_data(n_samples=100)
	X = (X - np.mean(X, axis=0)) / np.std(X, axis=0)  # 特征标准化
	W, b, losses = logistic_regression(X, Y, lr=0.1, epochs=1000, lambda_=0.5)
	print(f"最终损失: {losses[-1]}")
	plt.plot(losses)
	plt.title("Loss Curve")
	plt.xlabel("Epochs")
	plt.ylabel("Loss")
	plt.show()
	# 绘制决策边界
	plot_decision_boundary(X, Y, W, b)
		\end{lstlisting}
	\end{prob}
	
	\begin{sol}
解：\\
补充部分的代码如下：
\begin{lstlisting}[language=python,style=mystyle]
loss = -(1 / m) * np.sum(Y * np.log(A) + (1 - Y) * np.log(1 - A))
loss += (lambda_ / (2m)) * np.sum(W ** 2) ## 除以2是方便梯度计算, 注意我这里已经除以2了，后面梯度就不用除以2了,根据老师所讲这里需要除以m

......

dZ = A - Y  # 损失对预测的梯度
dW = (1 / m) * np.dot(X.T, dZ) + lambda_ * W /m  # 权重参数
db = (1 / m) * np.sum(dZ)  # 偏置
\end{lstlisting}

关键步骤如下：

在逻辑回归中，我们通过最小化损失函数来优化模型参数。损失函数由两部分组成：交叉熵损失和 L2 正则化项。目标是计算损失函数关于模型参数（权重 \(W\) 和偏置 \(b\)）的梯度。

\textbf{损失函数} \\
逻辑回归的损失函数由交叉熵损失和 L2 正则化项组成：
\[
\mathcal{L}(W, b) = -\frac{1}{m} \sum_{i=1}^{m} \left[ Y^{(i)} \log(A^{(i)}) + (1 - Y^{(i)}) \log(1 - A^{(i)}) \right] + \frac{\lambda}{2m} \sum_{i=1}^{n} W_i^2
\]


为了通过梯度下降方法优化模型，我们需要计算损失函数对于模型参数 \(W\) 和 \(b\) 的梯度。

\textbf{(a) 对权重 \(W\) 的梯度} \\
首先，我们需要对损失函数关于权重 \(W\) 求梯度。通过链式法则，我们可以得到：

\[
\frac{\partial \mathcal{L}}{\partial W} = \frac{\partial \mathcal{L}}{\partial Z} \cdot \frac{\partial Z}{\partial W}
\]
其中，\(\frac{\partial \mathcal{L}}{\partial Z} = A - Y\) 是交叉熵损失对于 \(Z\) 的梯度，这一步也是由链式法则直接计算得到的，且：
\[
\frac{\partial Z}{\partial W} = X^T
\]
因此，损失函数对于 \(W\) 的梯度为：
\[
\frac{\partial \mathcal{L}}{\partial W} = \frac{1}{m} \sum_{i=1}^{m} (A^{(i)} - Y^{(i)}) X^{(i)} + \lambda W/m
\]
其中，\( (A^{(i)} - Y^{(i)}) \) 是交叉熵损失对预测值 \(A^{(i)}\) 的梯度，\(\lambda W\) 是 L2 正则化项对 \(W\) 的梯度。
\textbf{(b) 对偏置 \(b\) 的梯度} \\
损失函数对于偏置 \(b\) 的梯度可以通过类似的方法计算：

\[
\frac{\partial \mathcal{L}}{\partial b} = \frac{\partial \mathcal{L}}{\partial Z} \cdot \frac{\partial Z}{\partial b}
\]
其中，\(\frac{\partial Z}{\partial b} = 1\)，所以损失函数对 \(b\) 的梯度为：
\[
\frac{\partial \mathcal{L}}{\partial b} = \frac{1}{m} \sum_{i=1}^{m} (A^{(i)} - Y^{(i)})
\]



在Colab平台运行，最终结果损失为：0.04380993306764013。
结果见图~\ref{fig:1}。
\begin{figure}[h!]
	\centering
	
\includegraphics[width=0.4\linewidth]{Figures/1}
	\caption{训练结果图}
	\label{fig:1}
\end{figure}

	\end{sol}
	
	\newpage
	
	\begin{prob}
		\textbf{L1 正则化和 L2 正则化的主要区别是什么？它们分别适用于什么场景？}
	\end{prob}
	
	\begin{sol}
		$$
		\begin{array}{|c|c|}
			\hline
			\textbf{特征} & \textbf{L1 正则化} \\
			\hline
			\textbf{公式} & \lambda \sum_{i=1}^{n} |w_i| \\
			\hline
			\textbf{惩罚方式} & \textbf{权重的绝对值（L1 范数）} \\
			\hline
			\textbf{对权重的影响} & \textbf{产生稀疏解（部分权重趋近于零）} \\
			\hline
			\textbf{模型稀疏性} & \textbf{更倾向于产生稀疏模型（有些权重为零）} \\
			\hline
			\textbf{优化算法} & \textbf{产生不连续的梯度，可能导致优化不稳定} \\
			\hline
		\end{array}
		$$
		\hspace{1cm} % 添加水平间距
		$$
		\begin{array}{|c|c|}
			\hline
			\textbf{特征} & \textbf{L2 正则化} \\
			\hline
			\textbf{公式} & \lambda \sum_{i=1}^{n} w_i^2 \\
			\hline
			\textbf{惩罚方式} & \textbf{权重的平方（L2 范数）} \\
			\hline
			\textbf{对权重的影响} & \textbf{使权重较小，但不完全为零} \\
			\hline
			\textbf{模型稀疏性} & \textbf{不倾向于产生稀疏模型，所有权重较小} \\
			\hline
			\textbf{优化算法} & \textbf{产生平滑的梯度，优化更稳定} \\
			\hline
		\end{array}
		$$
		
	L1正则化适用于高维数据，用于特征的选择，或者需要稀疏模型时；L2正则化适用于特征之间存在多重共线性时，避免过拟合，需要模型稳定训练时。
	
\begin{figure}[h!]
	\centering
	
\includegraphics[width=0.5\linewidth]{Figures/2}
	\caption{L1为什么稀疏}
	\label{fig:2}
\end{figure}
	对于L1稀疏的理解可以用一张图来展示，如图~\ref{fig:2}。
	来源：\href{https://freemind.pluskid.org/machine-learning/sparsity-and-some-basics-of-l1-regularization/}{Sparsity and Some Basics of L1 Regularization}.
	\end{sol}
	
	\newpage
	
	\begin{prob}
		\textbf{AdaGrad 和 RMSProp 的主要区别是什么？ Adam优化器结合了哪些优化算法的优点？它的主要超参数有哪些？}
	\end{prob}
	
	\begin{sol}
		解：\\
		\begin{enumerate}
			\item \textbf{主要区别：}\\AdaGrad 使用累计梯度的平方和用来控制学习率大小，这会导致学习率一直减小，一般凸优化问题中表现较好效果，而RMSProp使用指数加权移动平均(EMA) 计算梯度的历史平方和，使得梯度平方的历史信息逐渐淡化，使算法更适合处理非平稳目标,往往更适合深度学习（以及大多数非凸问题）,两者主要区别在代码中就体现为一行，见图~\ref{fig:3}
			\item \textbf{结合哪些算法的优点：}\\Adam优化器结合了RMSProp的自适应学习率,避免某些参数的梯度过大或过小，从而更稳定地更新参数的优点和Momentum的动量加速,避免梯度波动过大，有助于更快地找到最优解的优点
			\item \textbf{Adam的主要超参数：}\\有$\eta $学习率 $\beta_1$一阶动量衰减系数$\beta_2$二阶动量衰减系数, 避免除$0$的小量$\epsilon=1e-7$.
		\end{enumerate}
\begin{figure}[h!]
	\centering
	\includegraphics[width=1\linewidth]{Figures/3}
	\caption{AdaGrad 和 RMSProp 的主要区别}
	\label{fig:3}
\end{figure}
	\end{sol}
	
	\newpage
	
	\begin{prob}
		鞍点和局部最优值的区别是什么？如何通过 Hessian 矩阵的特征值来区分它们？
	\end{prob}
	
	\begin{sol}
	解：\\
\begin{enumerate}
	\item \textbf{区别：}\\
	局部最优值是指在某点$x^*$附近存在一个小的邻域，在该区域内的所有点的函数值$f(x)$均不小于(或不大于) $f(x^*)$，则$f(x^*)$称为局部最优值。这些点梯度为0.
	
	鞍点是指函数在某点$x_{saddle}$的梯度为零或者很小，但该点既不是局部极大值也不是局部极小值。换句话说，函数值在该点沿某些方向是极小值，而沿某些方向是极大值。这样的点很像是马鞍的形状
	
	\item \textbf{如何区分：}\\
	局部最优值的Hessian矩阵为正定或者负定矩阵(特征值都正或者都负)，鞍点的Hessian矩阵的特征值有正有负，如果Hessian矩阵为半正定或者半负定矩阵，则需要进一步判断。
\end{enumerate}
	\end{sol}

\end{document}
